% =====================================================================
%  TKCE — Methodology & Experiments  (AAAI draft)
% ---------------------------------------------------------------------
%  HOW TO USE THIS FILE
%  1. It compiles standalone (pdflatex) with the preamble below so you can
%     see the draft immediately.
%  2. To drop into the AAAI template: delete the preamble (between the
%     %%% PREAMBLE markers) and paste the two \section blocks into your
%     aaai25.tex. The AAAI style already provides amsmath/booktabs/graphicx.
%
%  IMAGE PLACEHOLDERS
%  Every figure has a grey box beginning "IMAGE PROMPT:". Copy that text
%  into your image-generation model, generate the PNG, save it under
%  paper/figures/<name>.png, then replace the placeholder box with the
%  commented-out \includegraphics line right below it.
%  Figures marked "GENERATED BY CODE" are produced by make_report.py
%  (e.g. cd_diagram.png) — just \includegraphics those directly.
% =====================================================================

%%% PREAMBLE (remove when pasting into the AAAI template) %%%%%%%%%%%%%%
\documentclass[11pt]{article}
\usepackage[margin=1in]{geometry}
\usepackage{amsmath,amssymb}
\usepackage{booktabs}
\usepackage{graphicx}
\usepackage{xcolor}
\usepackage{algorithm}
\usepackage{algpseudocode}
\usepackage{natbib}        % \citet / \citep (AAAI template already provides this)
\usepackage[hidelinks]{hyperref}
\usepackage{multirow}
\newcommand{\R}{\mathbb{R}}
\newcommand{\ind}{\mathbf{1}}
% Grey placeholder box that visibly shows the image prompt in the draft.
\newcommand{\imgprompt}[2]{%
  \begin{center}\fcolorbox{black}{gray!12}{%
  \begin{minipage}{0.92\linewidth}\small\vspace{3pt}
  \textbf{[IMAGE PLACEHOLDER — #1]}\\[2pt]
  \textbf{IMAGE PROMPT:} #2\vspace{3pt}
  \end{minipage}}\end{center}}
\begin{document}
%%% END PREAMBLE %%%%%%%%%%%%%%%%%%%%%%%%%%%%%%%%%%%%%%%%%%%%%%%%%%%%%%%%

% =====================================================================
\section{Methodology}
% =====================================================================

\subsection{Overview and Motivation}

Gradient-boosted decision trees (GBDTs) remain the state of the art on
medium-sized tabular data, and recent analyses attribute their edge over neural
networks (NNs) to three properties of tree partitions: robustness to
uninformative features, sensitivity to the axis-aligned orientation of the data
(NNs are rotationally invariant when they should not be), and the ability to fit
irregular, non-smooth target functions \cite{grinsztajn2022}. Our goal is to
\emph{transfer this tree inductive bias into an NN's representation}. We do so by
letting a tree ensemble define a similarity kernel over rows, learning an encoder
$\phi$ whose geometry reproduces that kernel through a contrastive objective, and
then predicting with a standard neural head (MLP or TabResNet) on top of $\phi$.
We call the method \textbf{TKCE} (Tree-Kernel Contrastive Embeddings).

\imgprompt{Figure 1 — Method overview}{A clean, professional academic
schematic (horizontal, left-to-right, flat vector style, muted blue/grey
palette, white background, no photorealism). Left: a small table of tabular
rows. Arrow to a box labelled ``Tree ensemble (XGBoost / Random Forest /
Mondrian)'' drawn as three small decision trees. Arrow to a matrix labelled
``Leaf co-occupancy kernel K(i,j)'' shown as a heatmap grid. Arrow to two
identical neural-network encoder blocks sharing weights (a Siamese pair)
labelled ``$\phi$'', with a ``contrastive loss'' symbol pulling similar points
together and pushing dissimilar points apart on a small sphere. From the encoder
output, two branches: top branch ``Two-stage: freeze $\phi$, train head''; bottom
branch ``Joint: train $\phi$ + head together''. Each branch ends in a box
``MLP / TabResNet head $\rightarrow$ prediction''. Label everything with small
clean sans-serif text.}

% \includegraphics[width=\linewidth]{figures/method_overview.png}

\subsection{Problem Setup and Notation}

We are given a tabular dataset $\mathcal{D}=\{(x_i,y_i)\}_{i=1}^{n}$ with
$x_i\in\R^{p}$ after preprocessing (numerical features standardized; categorical
features ordinal-encoded for trees, and either one-hot or learned embeddings for
NNs) and a target $y_i$ that is a class label (classification) or a real value
(regression). We learn an encoder $\phi_\theta:\R^{p}\!\to\!\mathcal{S}^{d-1}$
that maps a row to a unit-norm embedding, and a task head
$g_\psi$ that maps an embedding to a prediction.

\subsection{The Tree Kernel: What Is Similar}

A decision tree $t$ sends every row to exactly one leaf; write $\ell_t(x)$ for
that leaf index. Encoding ``which leaf'' as a one-hot vector $z_t(x)$ makes leaf
agreement a dot product, $\langle z_t(x_i),z_t(x_j)\rangle=\ind[\ell_t(x_i)=
\ell_t(x_j)]$. Averaging over an ensemble of $T$ trees gives the
\emph{leaf co-occupancy kernel}
\begin{equation}
K(x_i,x_j)\;=\;\frac{1}{T}\sum_{t=1}^{T}\ind\!\left[\ell_t(x_i)=\ell_t(x_j)\right]
\;\in\;[0,1].
\label{eq:kernel}
\end{equation}
$K$ is $1$ when two rows share a leaf in every tree (very similar) and $0$ when
they never co-occur. This is not an arbitrary choice: random tree partitions
(a Mondrian process) give a random-feature approximation to the isotropic Laplace
kernel $k(x,x')=\exp(-\lambda\lVert x-x'\rVert_1)$ \cite{balog2016}, so
Eq.~\eqref{eq:kernel} is a fast, data-adaptive stand-in for a smooth classical
similarity that decays with distance.

\imgprompt{Figure 2 — Tree kernel intuition}{A minimalist explanatory diagram,
flat vector style, white background. Show a 2D feature space partitioned by
axis-aligned rectangles (like a decision-tree partition) three times side by side
(``Tree 1'', ``Tree 2'', ``Tree 3''), each with slightly different rectangles.
In each panel place the same three points A, B, C. A and B fall in the same
rectangle in 3 of 3 trees (draw them close); C falls in a different rectangle
(draw it apart). To the right, a small formula callout: ``K(A,B)=1.0 (same leaf
everywhere), K(A,C)=0.33''. Caption-free, clean labels, soft colours to shade
the leaf regions.}

% \includegraphics[width=\linewidth]{figures/kernel_intuition.png}

We consider three kernel \emph{sources}, all plugged into Eq.~\eqref{eq:kernel}:
\begin{itemize}
  \item \textbf{GBT (target-led).} Leaves come from a gradient-boosted forest fit
  on $y$. The similarity is \emph{task-relevant}: it reflects the structure the
  trees found useful for prediction. This is our primary kernel.
  \item \textbf{Random Forest.} Bagged, supervised, but less greedy than boosting.
  \item \textbf{Mondrian (label-free).} Axis-aligned partitions drawn at random
  \emph{without} using $y$ \cite{balog2016}. This yields a purely unsupervised,
  task-agnostic prior and lets us ask whether even a label-free tree structure
  helps.
\end{itemize}

\subsection{Contrastive Encoder: Turning Similarity into Geometry}

We train a Siamese encoder $\phi$ (one network, shared weights, applied to both
members of a pair) so that embedding geometry reproduces $K$. Because $\phi$ is
$L_2$-normalized, $\langle\phi(x_i),\phi(x_j)\rangle$ is a cosine similarity. We
use two objectives.

\paragraph{InfoNCE.} From pairs $(x_i,x_i^+)$ that are \emph{positive} under the
kernel (i.e.\ $K(x_i,x_i^+)\!\ge\!\tau_{\text{pos}}$), with the other rows in the
batch as negatives and temperature $\tau$,
\begin{equation}
\mathcal{L}_{\text{NCE}}=-\log\frac{\exp(\langle\phi(x_i),\phi(x_i^+)\rangle/\tau)}
{\sum_{k}\exp(\langle\phi(x_i),\phi(x_k)\rangle/\tau)}.
\label{eq:infonce}
\end{equation}

\paragraph{Kernel regression.} A soft alternative that matches the dot product to
the exact kernel value,
$\mathcal{L}_{\text{KR}}=\big(\langle\phi(x_i),\phi(x_j)\rangle-K(x_i,x_j)\big)^2.$

Positive pairs are mined by sampling candidate pairs and keeping those with
$K\!\ge\!\tau_{\text{pos}}$, which avoids materializing the $O(n^2)$ kernel and
scales to large $n$.

\paragraph{Loss family.} InfoNCE and kernel regression are our defaults, but the
choice of contrastive objective is itself an experimental variable. We evaluate
\textbf{seven} losses spanning the standard families: InfoNCE (NT-Xent),
kernel regression, the margin \emph{contrastive} pair loss \cite{hadsell2006},
\emph{triplet} loss \cite{schroff2015}, \emph{supervised contrastive}
(multiple positives per anchor, a natural fit since the kernel yields many
positives) \cite{khosla2020}, \emph{anisotropic} InfoNCE \cite{aninfonce}, and a
\emph{CLIP}-style symmetric InfoNCE with a learnable temperature \cite{clip}.
All but supervised-contrastive and kernel regression share the in-batch
anchor/positive path; the ablation is in Sec.~``Loss-function ablation''.

\subsection{Two Regimes for Injecting the Tree Bias}

A central question of this work is \emph{how} the tree bias should enter the
network. We study two regimes (Fig.~3).

\paragraph{(A) Two-stage (decoupled).} First pretrain $\phi$ with the contrastive
loss (no labels used in the pretraining objective beyond those already baked into
a supervised kernel). Then train a head $g$ on the frozen (or fine-tuned)
embeddings with the task loss $\mathcal{L}_{\text{task}}$ (cross-entropy or MSE).
This is the standard self-supervised ``pretrain then probe'' protocol.

\paragraph{(B) Joint (end-to-end).} Train encoder and head together, adding the
contrastive term as an auxiliary regularizer, analogous to learning an embedding
layer jointly with the task:
\begin{equation}
\mathcal{L}_{\text{joint}}=\mathcal{L}_{\text{task}}\big(g(\phi(x)),y\big)
\;+\;\lambda\,\mathcal{L}_{\text{con}}(\phi;K).
\label{eq:joint}
\end{equation}
Here $\lambda\!\to\!0$ recovers a plain NN and large $\lambda$ recovers pure
embedding learning; $\lambda$ trades off task fit against fidelity to the tree
similarity.

\imgprompt{Figure 3 — Two-stage vs.\ joint}{A side-by-side comparison diagram,
flat vector style, white background, two panels. LEFT panel titled ``(A)
Two-stage'': a frozen (padlock icon) encoder $\phi$ feeding a head; below it a
timeline ``Stage 1: contrastive pretrain $\rightarrow$ Stage 2: train head''.
RIGHT panel titled ``(B) Joint'': encoder $\phi$ and head in one block with two
loss arrows entering it, one labelled ``task loss'' and one labelled
``$\lambda\times$ contrastive loss'', trained simultaneously (circular-arrow
icon). Keep icons simple and labels in clean sans-serif.}

% \includegraphics[width=\linewidth]{figures/two_regimes.png}

\subsection{Architectures and the Fixed-Head Protocol}
The encoder $\phi$ is a small MLP with batch-norm and $L_2$-normalized output.
Heads are (i) an \textbf{MLP} and (ii) a \textbf{TabResNet} (pre-norm residual
blocks) \cite{gorishniy2021}. Algorithm~\ref{alg:tkce} summarizes both regimes.

\paragraph{Controlled comparison via a fixed head.} To isolate the effect of the
\emph{representation}, we hold the downstream head \emph{fixed} (identical hidden
sizes and dropout for the MLP head, identical depth/width for the TabResNet head)
across every method that reads a representation --- the raw-feature NN, the
leaf-one-hot NN, and all TKCE variants. Only the \emph{encoder} $\phi$ (its
width, depth, dropout, embedding dimension, learning rate) and the kernel/
contrastive hyperparameters are tuned. Thus any difference between, say, an MLP
on raw features and an MLP on TKCE embeddings is attributable to the input
representation rather than to head capacity. The strong deep baselines
(FT-Transformer, MLP+numerical-embeddings) are integrated architectures rather
than a head on a fixed representation, so they retain their own tuned designs
and serve as the competitive upper bar.

\begin{algorithm}[t]
\caption{TKCE training (both regimes)}
\label{alg:tkce}
\begin{algorithmic}[1]
\Require data $\mathcal{D}$, kernel source $s\in\{$gbt, rf, mondrian$\}$,
regime $\in\{$two-stage, joint$\}$, weight $\lambda$
\State Fit tree ensemble of source $s$ on train; extract leaves $\ell_t(\cdot)$
\State Define $K$ by Eq.~\eqref{eq:kernel}; mine positive pairs
\If{regime = two-stage}
  \State Pretrain $\phi$ with $\mathcal{L}_{\text{NCE}}$ (Eq.~\ref{eq:infonce})
  \State Freeze/fine-tune $\phi$; train head $g$ on $\phi(x)$ with
  $\mathcal{L}_{\text{task}}$
\Else \Comment{joint}
  \State Train $\{\phi,g\}$ with $\mathcal{L}_{\text{joint}}$
  (Eq.~\ref{eq:joint})
\EndIf
\State \Return $g\circ\phi$
\end{algorithmic}
\end{algorithm}

% =====================================================================
\section{Experiments}
% =====================================================================

\subsection{Research Questions}
\begin{description}
  \item[RQ1] Does TKCE close the tree$\rightarrow$NN performance gap relative to
  NNs trained on raw features?
  \item[RQ2] Which injection regime is better — two-stage or joint?
  \item[RQ3] How much does the kernel \emph{source} matter (supervised GBT vs.\
  label-free Mondrian)?
  \item[RQ4] Does the learned contrastive embedding beat the cheap alternative of
  feeding raw tree-leaf indices to an NN, and does it rival strong tabular-DL
  baselines?
  \item[RQ5] Does \emph{contrastive} representation learning beat a \emph{linear}
  PCA embedding of the same tree features, through the same fixed head?
\end{description}

\subsection{Datasets}
We use the tabular benchmark of \citet{grinsztajn2022}, curated specifically to
compare trees and NNs. It spans four quadrants --- numerical vs.\ heterogeneous
(numerical$+$categorical) features $\times$ classification vs.\ regression ---
accessed by OpenML benchmark-suite id (337, 336, 334, 335 respectively). Each
dataset is split chronology-free into train/validation/test; numerical features
are standardized on train only, and regression targets are standardized (metrics
reported on the original scale).

\subsection{Baselines and Model Grid}
Table~\ref{tab:models} lists all models. Crucially, we include (i) strong tree
baselines tuned as hard as the NNs, (ii) strong \emph{deep} baselines
(FT-Transformer \cite{gorishniy2021}; an MLP with periodic numerical embeddings
\cite{gorishniy2022}), (iii) the classic \emph{leaf-one-hot$\to$MLP}
baseline \cite{he2014}, i.e.\ feeding raw GBDT leaf indices as one-hot features
--- the cheap alternative our contrastive embedding must beat (RQ4) --- and
(iv) a \emph{PCA} baseline that replaces the contrastive encoder with a
\emph{linear} projection of the same leaf features. Because the leaf one-hot is
the tree kernel's explicit feature map, this PCA is the linear spectral embedding
of the tree kernel; comparing it to TKCE through the \emph{same} fixed head
isolates the value of \emph{contrastive} (non-linear) representation learning
over plain dimensionality reduction (RQ5).

\begin{table}[t]\centering\small
\caption{Model grid. Every model is tuned with an equal Optuna budget.}
\label{tab:models}
\begin{tabular}{ll}
\toprule
Group & Models \\
\midrule
Tree baselines      & XGBoost, LightGBM, CatBoost, Random Forest \\
NN on raw features  & MLP, TabResNet \\
Strong deep tabular & FT-Transformer, MLP+numerical-embed. \\
Tree-feature baselines & leaf-one-hot $\to$ MLP; \textbf{PCA of leaf features} $\to$ MLP/TabResNet \\
\textbf{TKCE two-stage} & MLP / TabResNet on $\phi$ (GBT, Mondrian) \\
\textbf{TKCE joint}     & MLP / TabResNet $+$ contrastive (GBT) \\
\bottomrule
\end{tabular}
\end{table}

\subsection{Protocol}
Following \citet{grinsztajn2022}, every model with tunable hyperparameters
receives an \emph{equal budget} of 100 Optuna trials maximizing the validation
metric, per model per dataset --- this covers the trees, the strong deep
baselines, and the TKCE \emph{encoder}/kernel. The fixed-head controls
(raw-feature MLP/TabResNet) have no hyperparameters to search and are trained
once with the standard head configuration. The winning configuration is then
evaluated over multiple random seeds; we report mean $\pm$ std. The primary
metric is ROC-AUC (classification) and $R^2$ (regression). To aggregate across
datasets we compute per-dataset ranks, average ranks, a critical-difference
(Nemenyi) diagram \cite{demsar2006}, and Wilcoxon signed-rank tests against a
reference model.

\subsection{Main Results (RQ1, RQ4)}
Table~\ref{tab:main} reports the primary metric per dataset. Table~\ref{tab:gap}
quantifies the \emph{gap closed},
$\big(\text{best}_{\text{TKCE}}-\text{best}_{\text{raw NN}}\big)/
\big(\text{best}_{\text{tree}}-\text{best}_{\text{raw NN}}\big)$, on datasets
where trees lead. Figure~4 shows the aggregate ranking.

\begin{table}[t]\centering\scriptsize
\caption{Primary metric (mean$\pm$std over seeds); best per row in bold.
\textbf{Auto-generated by \texttt{make\_report.py} $\to$ main\_results.csv;
paste the produced rows here.}}
\label{tab:main}
\begin{tabular}{llccc}
\toprule
Dataset & Metric & Best tree & Best raw NN & \textbf{Best TKCE} \\
\midrule
$\langle$dataset$\rangle$ & AUC/$R^2$ & -- & -- & -- \\
\multicolumn{5}{c}{\emph{(rows filled from main\_results.csv)}} \\
\bottomrule
\end{tabular}
\end{table}

\begin{table}[t]\centering\small
\caption{Gap analysis (from \texttt{gap\_analysis.csv}). A positive
\emph{gap closed} means TKCE moves the NN toward the tree ceiling.}
\label{tab:gap}
\begin{tabular}{lcccc}
\toprule
Dataset & best tree & best raw NN & best TKCE & gap closed \\
\midrule
$\langle$dataset$\rangle$ & -- & -- & -- & --\,\% \\
\bottomrule
\end{tabular}
\end{table}

\paragraph{Figure 4 — Aggregate ranking (GENERATED BY CODE).}
Use \texttt{cd\_diagram.png} produced by \texttt{make\_report.py} directly.
% \includegraphics[width=\linewidth]{figures/cd_diagram.png}

\imgprompt{Figure 5 — Gap-closed bar chart}{A clean academic bar chart, flat
vector style, white background, colour-blind-safe palette. X-axis: dataset names
(rotated). Y-axis: ``\% of tree--NN gap closed by TKCE'', range 0--100\%. One bar
per dataset, a dashed horizontal line at the median. Title ``TKCE closes the
tree--NN gap''. Minimal gridlines, no 3D. (You can also regenerate this from
gap\_analysis.csv.)}

% \includegraphics[width=\linewidth]{figures/gap_closed.png}

\subsection{Two-Stage vs.\ Joint (RQ2)}
Table (from \texttt{twostage\_vs\_joint.csv}) and Figure~6 compare the regimes.
Our hypothesis, motivated by preliminary runs, is that the \emph{joint} regime
closes more of the gap on tree-dominant datasets, while the frozen two-stage
embedding can lag when the task signal is subtle.

\imgprompt{Figure 6 — Regime comparison}{A grouped bar chart, flat vector style,
white background. X-axis: dataset names. For each dataset two adjacent bars:
``two-stage'' and ``joint'', y-axis = primary metric. A short annotation arrow
highlighting datasets where joint clearly wins. Clean, publication-quality,
colour-blind-safe.}

% \includegraphics[width=\linewidth]{figures/twostage_vs_joint.png}

For the joint regime we sweep the contrastive weight $\lambda$
(Eq.~\ref{eq:joint}); Fig.~\ref{fig:lambda} shows performance across $\lambda$,
with $\lambda\!=\!0$ (no tree bias) and the tree/raw-NN references. A peak above
$\lambda\!=\!0$ confirms the contrastive term helps and locates the sweet spot.

% Figure is GENERATED BY CODE (run_lambda_sweep.py).
\begin{figure}[t]\centering
\includegraphics[width=0.9\linewidth]{figures/lambda_sweep.png}
\caption{Effect of the contrastive weight $\lambda$ in the joint regime
(generated by \texttt{run\_lambda\_sweep.py}). $\lambda\!=\!0$ is a plain NN;
dashed/dotted lines are the raw-MLP and best-tree references.}
\label{fig:lambda}
\end{figure}

\subsection{Kernel-Source Ablation (RQ3)}
Table (from \texttt{kernel\_ablation.csv}) compares the supervised GBT kernel
against the label-free Mondrian kernel (both with an MLP head). We expect the
target-led GBT kernel to lead on task performance, with the Mondrian result
indicating how much benefit survives \emph{without} any label information.

\subsection{Loss-Function Ablation}
Holding the kernel and the fixed head constant, we sweep the seven contrastive
losses and measure the downstream metric of the resulting two-stage embedding
(Fig.~\ref{fig:lossabl}). Raw-MLP and best-tree references mark the useful range:
a loss earns its place only if its embedding beats raw features.

% Figure is GENERATED BY CODE (run_loss_ablation.py).
\begin{figure}[t]\centering
\includegraphics[width=\linewidth]{figures/loss_ablation.png}
\caption{Contrastive-loss ablation (generated by \texttt{run\_loss\_ablation.py}):
downstream metric of a two-stage TKCE embedding trained with each loss, with
raw-MLP (dashed) and best-tree (dotted) references.}
\label{fig:lossabl}
\end{figure}

\subsection{Contrastive vs.\ Linear (PCA) Embedding (RQ5)}
To test whether the \emph{contrastive} objective is necessary --- rather than any
low-dimensional projection of the tree features --- we compare TKCE against a PCA
embedding of the same leaf-one-hot features, fed to the \emph{same} fixed head
and tuned over the same number of components. Since the leaf one-hot is the
kernel's explicit feature map, PCA here is the linear spectral embedding of the
tree kernel. If TKCE beats PCA, the gain is attributable to \emph{non-linear
contrastive} learning, not merely to compressing the tree features; if not, a
linear projection suffices and the encoder is unnecessary. Results appear in the
main table (the \texttt{pca\_gbt\_*} rows) and the aggregate ranking.

\subsection{Why It Works: Mechanism (Analysis)}
To connect results to the three tree advantages \cite{grinsztajn2022}, we run
controlled probes: (i) injecting $k$ uninformative noise features and measuring
robustness; (ii) applying a random rotation to the features (which should hurt
axis-aligned methods) and measuring degradation; (iii) targets of increasing
irregularity. We report where TKCE tracks the trees and where it does not.

% Figure 7 is GENERATED BY CODE (run_mechanism.py -> figures/mechanism.png).
\begin{figure*}[t]\centering
\includegraphics[width=\linewidth]{figures/mechanism.png}
\caption{Mechanism probes (generated by \texttt{run\_mechanism.py}). (a)
Robustness as uninformative features are added; (b) degradation under feature
rotation --- axis-aligned methods (trees, and TKCE if it inherited the bias)
should drop while a rotation-invariant MLP stays flat; (c) increasingly irregular
(high-frequency) targets. We report where TKCE tracks the trees.}
\label{fig:mechanism}
\end{figure*}

\subsection{Data Efficiency (Analysis)}
Trees excel in the low-data regime. Figure~8 plots performance vs.\ training-set
fraction; we test whether TKCE inherits this low-data strength.

% Figure 8 is GENERATED BY CODE (run_data_efficiency.py).
\begin{figure}[t]\centering
\includegraphics[width=\linewidth]{figures/data_efficiency.png}
\caption{Data-efficiency (generated by \texttt{run\_data\_efficiency.py}):
primary metric vs.\ training-set fraction, with $\pm$std bands over seeds. We
test whether TKCE sits above the raw NN especially at small fractions.}
\label{fig:dataeff}
\end{figure}

\subsection{Cost}
We report per-model training cost (wall-clock, parameters). Trees train in
seconds; TKCE adds a forest fit plus contrastive pretraining. We discuss whether
the accuracy gain justifies the added cost.

% ---------------------------------------------------------------------
\begin{thebibliography}{9}\small
\bibitem{grinsztajn2022} L.~Grinsztajn, E.~Oyallon, G.~Varoquaux.
Why do tree-based models still outperform deep learning on typical tabular data?
\emph{NeurIPS Datasets \& Benchmarks}, 2022. arXiv:2207.08815.
\bibitem{gorishniy2021} Y.~Gorishniy et al. Revisiting Deep Learning Models for
Tabular Data. \emph{NeurIPS}, 2021. arXiv:2106.11959.
\bibitem{gorishniy2022} Y.~Gorishniy, I.~Rubachev, A.~Babenko. On Embeddings for
Numerical Features in Tabular Deep Learning. \emph{NeurIPS}, 2022.
arXiv:2203.05556.
\bibitem{balog2016} M.~Balog et al. The Mondrian Kernel. \emph{UAI}, 2016.
arXiv:1606.05241.
\bibitem{he2014} X.~He et al. Practical Lessons from Predicting Clicks on Ads at
Facebook. \emph{ADKDD}, 2014.
\bibitem{demsar2006} J.~Dem\v{s}ar. Statistical Comparisons of Classifiers over
Multiple Data Sets. \emph{JMLR}, 2006.
\bibitem{hadsell2006} R.~Hadsell, S.~Chopra, Y.~LeCun. Dimensionality Reduction
by Learning an Invariant Mapping. \emph{CVPR}, 2006.
\bibitem{schroff2015} F.~Schroff, D.~Kalenichenko, J.~Philbin. FaceNet: A Unified
Embedding for Face Recognition and Clustering. \emph{CVPR}, 2015.
arXiv:1503.03832.
\bibitem{khosla2020} P.~Khosla et al. Supervised Contrastive Learning.
\emph{NeurIPS}, 2020. arXiv:2004.11362.
\bibitem{aninfonce} A.~Rusak et al. AnInfoNCE: Identifying the Latent Factors in
Anisotropic Contrastive Learning. 2024. arXiv:2407.00143.
\bibitem{clip} A.~Radford et al. Learning Transferable Visual Models From Natural
Language Supervision (CLIP). \emph{ICML}, 2021. arXiv:2103.00020.
\end{thebibliography}

\end{document}
